\section{Relation to Note}

This section clarifies the relation between this paper and the corresponding note-style exposition. The aim is not to summarize the theory, but to explain how the same structure appears in different descriptive layers.

\subsection{Difference in Descriptive Layer}

This paper is written as an L3 formal layer. The note essay is written as an L1 intuitive layer. They concern the same target, but their purposes differ. The paper fixes definitions, equations, and relations. The note essay supports intuitive understanding through accessible language and metaphor.

The relevant relation is therefore structural correspondence, not literal identity of wording.

\subsection{Correspondences}

The concepts used in this paper correspond to intuitive expressions in the note essay:
\begin{align}
  \I &\leftrightarrow \text{local concentration, bias, or structured region},\\
  \J &\leftrightarrow \text{flow, repetition, or ease of passage},\\
  \PhiI &\leftrightarrow \text{return tendency or stable place},\\
  \text{attractor} &\leftrightarrow \text{returning to the same region or preserving form}.
\end{align}
These are not direct translations. They are expressions of the same phenomenon at different resolution.

\subsection{Development and Form}

The treatment of development and morphogenesis in Section 6 corresponds to the note-level idea that repeated changes form flows and that flows stabilize into form. Both descriptions avoid taking design or purpose as primitives.

\subsection{Definition of Life}

Section 7 defines life as sustained stable motion. In the note essay, this appears as the idea that life is not an object but a state. The two expressions refer to the same structure at different levels of precision.

\subsection{Probability Revision}

The probability revision in Section 8 corresponds to the note-level statement that life is not best understood as a miracle-like exact hit. The shared structure is the shift from a point target to a region, and from exact matching to arrival and persistence.

\subsection{Connection to Intelligence}

The treatment of intelligence as a second attractor in Section 9 corresponds to a note-level description of a flow that appears to think. The note essay does not specify the full structure, but it preserves the distinction between life and intelligence as different stable organizations.

\subsection{Relation to Appendices}

Intermediate appendices can bridge the note essay and this paper. They may introduce state, flow, and stability with minimal mathematics. Future versions can be updated to reflect the structure fixed here.

\subsection{Summary}

This paper and the note essay differ by descriptive form, not by subject. The paper fixes the formal structure. The note essay presents the same structure experientially. Each can be read independently, but the correspondence improves coherence across layers.
